\documentclass[10pt]{article}
\providecommand{\FIGDIR}{fig}
\newcommand{\DIGESTDATE}{19 August 2026}
% ===================== PM Research Watch - shared preamble =====================
\usepackage[T1]{fontenc}
\usepackage[utf8]{inputenc}
\usepackage{amsmath,amssymb}
\usepackage{tgpagella}
\usepackage{tgheros}
\renewcommand{\sfdefault}{qhv}
\renewcommand{\ttdefault}{lmtt}
\usepackage{microtype}
\usepackage[english]{babel}

\usepackage[a4paper,top=17mm,bottom=17mm,left=17mm,right=17mm,
            headsep=6mm,footskip=10mm]{geometry}
\usepackage{graphicx}
\usepackage{xcolor}
\usepackage{colortbl}
\usepackage{tikz}
\usetikzlibrary{calc}
\usepackage[most]{tcolorbox}
\usepackage{titlesec}
\usepackage{fancyhdr}
\usepackage{booktabs}
\usepackage{array}
\usepackage{longtable}
\usepackage{enumitem}
\usepackage{multicol}
\usepackage{ragged2e}
\usepackage[hidelinks]{hyperref}

% ------------------------------------------------------------------ palette
\definecolor{Ink}    {HTML}{14262E}
\definecolor{Deep}   {HTML}{0F4C5C}
\definecolor{Teal}   {HTML}{1F7A8C}
\definecolor{Sky}    {HTML}{4FA3C4}
\definecolor{Sage}   {HTML}{6FA287}
\definecolor{Amber}  {HTML}{E3A008}
\definecolor{Clay}   {HTML}{B5643C}
\definecolor{Coral}  {HTML}{D1495B}
\definecolor{Violet} {HTML}{6A5B9E}
\definecolor{Slate}  {HTML}{4A6572}
\definecolor{Paper}  {HTML}{FBF9F5}
\definecolor{Mist}   {HTML}{EEEAE2}
\definecolor{Rule}   {HTML}{D9D3C9}

\pagecolor{Paper}
\color{Ink}

\hypersetup{colorlinks=true, linkcolor=Deep, urlcolor=Teal, citecolor=Teal,
            pdftitle={PM Research Watch}, pdfauthor={Automated daily digest}}

% ------------------------------------------------------------------ headings
\titleformat{\section}
  {\sffamily\bfseries\large\color{Deep}}{}{0pt}
  {}[{\vspace{1.2mm}\color{Amber}\titlerule[1.1pt]}]
\titlespacing*{\section}{0pt}{5.5mm}{2.8mm}

\titleformat{\subsection}
  {\sffamily\bfseries\normalsize\color{Teal}}{}{0pt}{}
\titlespacing*{\subsection}{0pt}{3.6mm}{1.4mm}

\setlist[itemize]{leftmargin=4.2mm, itemsep=1.1mm, topsep=1.3mm, parsep=0pt,
                  label=\textcolor{Teal}{\footnotesize$\blacktriangleright$}}

% ------------------------------------------------------------------ header/footer
\pagestyle{fancy}
\fancyhf{}
\renewcommand{\headrulewidth}{0pt}
\renewcommand{\footrulewidth}{0pt}
\fancyhead[L]{\sffamily\scriptsize\color{Slate}PM RESEARCH WATCH\;\textcolor{Rule}{\textbar}\;\DIGESTDATE}
\fancyhead[R]{\sffamily\scriptsize\color{Slate}Particulate matter monitoring \& health effects}
\fancyfoot[C]{\sffamily\scriptsize\color{Slate}\thepage}

% ------------------------------------------------------------------ components
% metric tile
\newcommand{\metric}[3]{%
  \begin{tcolorbox}[enhanced, colback=white, colframe=Rule, boxrule=0.5pt,
    arc=1.6mm, left=2.2mm, right=2.2mm, top=1.8mm, bottom=1.8mm,
    borderline west={1.7mm}{0pt}{#3}, nobeforeafter, valign=center,
    height=17mm, width=\linewidth]
    \RaggedRight
    {\sffamily\bfseries\fontsize{17}{17}\selectfont\color{#3}#1}\\[0.5mm]
    {\sffamily\fontsize{6.5}{7.4}\selectfont\color{Slate}#2}
  \end{tcolorbox}}

% pull-quote / key insight box
\newtcolorbox{keybox}[1][]{enhanced, colback=Mist, colframe=Mist,
  arc=1.6mm, left=3.2mm, right=3.2mm, top=2.4mm, bottom=2.4mm,
  borderline west={1.7mm}{0pt}{Amber}, #1}

% subtopic band
\newcommand{\band}[2]{%
  \vspace{3.2mm}\noindent
  \begin{tcolorbox}[enhanced, colback=#2, colframe=#2, arc=1.2mm,
    left=2.6mm, right=2.6mm, top=1.1mm, bottom=1.1mm, boxrule=0pt,
    width=\linewidth, nobeforeafter]
    {\sffamily\bfseries\small\color{white}#1}
  \end{tcolorbox}\vspace{1.6mm}}

% one paper entry: {short cite}{title}{journal}{doi}{body}
\newcommand{\paperentry}[5]{%
  \par\vspace{1.9mm}\noindent
  {\sffamily\bfseries\footnotesize\color{Ink}#1}\;%
  {\sffamily\fontsize{7}{8}\selectfont\color{Slate}#3}\\[0.4mm]
  {\footnotesize\itshape\color{Deep}#2}\\[0.6mm]
  {\footnotesize #4}\\[0.3mm]
  {\fontsize{6.6}{7.6}\selectfont\color{Slate}\texttt{doi:}\,\href{https://doi.org/#5}{\texttt{\detokenize{#5}}}}
  \par}

% caption style
\newcommand{\figcap}[1]{\par\vspace{1.0mm}
  {\fontsize{7.2}{8.6}\selectfont\color{Slate}\textsf{\textbf{Fig.}}\; #1\par}\vspace{1.5mm}}

\providecommand{\permil}{\textperthousand}

\setlength{\parindent}{0pt}
\setlength{\parskip}{0pt}
\renewcommand{\arraystretch}{1.18}


\begin{document}

% ==================================================================== MASTHEAD
\thispagestyle{empty}
\begin{tikzpicture}[remember picture, overlay]
  \fill[Deep] (current page.north west) rectangle
      ($(current page.north east)+(0,-67mm)$);
  \fill[Amber] ($(current page.north west)+(0,-67mm)$) rectangle
      ($(current page.north east)+(0,-68.4mm)$);
  \foreach \x/\y/\r in {22/12/0.9, 41/26/1.5, 63/9/0.7, 88/31/1.1, 112/17/1.9,
                        138/29/0.8, 159/13/1.3, 178/54/1.0, 196/21/1.6,
                        30/44/0.6, 74/57/1.2, 125/48/0.9, 168/6/0.8, 191/59/0.7,
                        52/18/1.0, 100/60/0.6, 148/20/1.1, 66/49/1.4, 205/45/0.8}
    \fill[Sky, opacity=0.30] ($(current page.north west)+(\x mm,-\y mm)$) circle (\r mm);
\end{tikzpicture}

\vspace*{4mm}
{\sffamily\fontsize{7.6}{9}\selectfont\color{Sky}%
 AUTOMATED DAILY LITERATURE DIGEST\;\textbullet\;PhD RESEARCH WATCH\par}
\vspace{2.2mm}
{\sffamily\bfseries\fontsize{27}{29}\selectfont\color{white}%
 Particulate Matter\par}
{\sffamily\fontsize{27}{29}\selectfont\color{Sky}%
 Monitoring \& Health Effects\par}
\vspace{5.0mm}
{\sffamily\fontsize{8.4}{10}\selectfont\color{white}%
 Issue 2026-08-19\;\;\textcolor{Amber}{\textbar}\;\;Entry window 19 August 2026\;\;%
 \textcolor{Amber}{\textbar}\;\;22 records in scope, 13 rejected\par}
\vspace{18mm}

% -------------------------------------------------------------- metric strip
\noindent
\begin{minipage}[t]{0.190\linewidth}\metric{22}{RECORDS IN SCOPE\\(13 REJECTED OFF-TOPIC)}{Deep}\end{minipage}\hfill
\begin{minipage}[t]{0.190\linewidth}\metric{9}{SUBTOPIC CLUSTERS\\REPRESENTED}{Teal}\end{minipage}\hfill
\begin{minipage}[t]{0.190\linewidth}\metric{6}{ABSTRACT-FREE\\PUBLISHER DEPOSITS}{Sky}\end{minipage}\hfill
\begin{minipage}[t]{0.190\linewidth}\metric{5}{EFFECT ESTIMATES\\EXTRACTED}{Amber}\end{minipage}\hfill
\begin{minipage}[t]{0.190\linewidth}\metric{3}{REVIEWS /\\EVIDENCE SYNTHESES}{Clay}\end{minipage}

\vspace{6mm}

% -------------------------------------------------------------- signal of day
\section*{\sffamily\bfseries\color{Deep}Signal of the day}
\vspace{-1.5mm}{\color{Amber}\rule{\linewidth}{1.1pt}}\vspace{2.5mm}

\begin{keybox}
{\sffamily\bfseries\small\color{Deep}Cumulative exposure \emph{duration}, not the day's
concentration, is where today's hazard signal sits --- and the day's largest acute study
shows why.}\\[1.4mm]
\footnotesize
Sivakumar \& Kurian (\textit{Environ Toxicol}) exposed female Wistar rats to
250\,$\mu$g/m$^3$ PM$_{2.5}$, 3\,h/day, for 1, 7, 14 or 21 days at a fixed concentration.
At 1--7 days the cardiac mitochondrial response was \textbf{adaptive and fully reversible}:
transient oxidative stress and quality-control activation, with respiration, ETC complex
activity, ATP synthesis and ex vivo cardiac function all indistinguishable from control
after a 24\,h pollutant-free washout. From \textbf{14 days} the same daily dose produced
coordinated suppression of respiration and ATP synthesis, mtDNA depletion and
intramitochondrial metal accumulation --- and these deficits \textbf{did not recover within
the 24\,h washout window}. Concentration was constant throughout; only duration changed.
Set against this, Park et al.\ (\textit{Gut Liver}) ran the day's largest acute design
--- 29{,}167 Korean cirrhosis patients, time-stratified case-crossover, 1\,km daily
PM$_{2.5}$ --- and recovered \textbf{aOR 1.02 (95\% CI 1.00--1.04) per IQR}, a lower bound
sitting exactly on the null, with associations \emph{attenuating} as the averaging window
lengthened. Two designs, opposite time-scales, same message for instrumentation: a
day-resolved mass metric is near its resolution floor, while the cumulative-duration metric
that the animal work implicates is precisely what a continuously-logging personal or
distributed sensor can deliver and a regulatory site cannot.
\end{keybox}

\vspace{3mm}

% -------------------------------------------------------------- exec brief
\section*{\sffamily\bfseries\color{Deep}What changed today}
\vspace{-1.5mm}{\color{Amber}\rule{\linewidth}{1.1pt}}\vspace{2.5mm}

\footnotesize
\begin{itemize}
\item \textbf{Filtration works; the model that predicts it does not transfer.} Xia et al.\
measured 33 spaces in 12 child-serving facilities before and after enhanced HVAC filtration
and portable air cleaners. PM$_{2.5}$ fell 45--54\% in standard schools and the nursery
(mixed-effects estimate $\sim$51\% for the PACs alone), but barely moved in childcare
houses. Crucially, single-zone mass-balance prediction held for regular classrooms and
failed for portable buildings, very large rooms and houses repurposed for childcare. The
buildings where the physics is unreliable are exactly the ones that need instrumenting.

\item \textbf{Proteomics enters as an effect-modification layer, not a main-effect one.}
He et al.\ screened 2{,}254 plasma proteins against LUR-modelled PM$_{2.5}$, PM$_{\mathrm{coarse}}$,
NO$_2$ and NO in 51{,}428 participants and report protein\,$\times$\,pollutant
\emph{interactions} on epilepsy risk (CIT--PM$_{2.5}$, PTPRR--NO$_2$, BCL2L15--API), with
hippocampal-volume and white-matter-hyperintensity correlates. The multiplicity burden here
is roughly 9{,}000 interaction tests and the abstract reports no corrected significance
level or interaction magnitude --- treat as a candidate list, not an established mechanism.

\item \textbf{PM$_{2.5}$ mass is again shown to be a poor proxy for a co-located particle
property.} Yang et al.\ (Taiyuan, month-long winter campaign) found immersion-mode INP
concentration rose from 0.61 to 7.47\,L$^{-1}$ (95\% CI 6.64--8.41) during a desert-dust
intrusion, yet showed $|r|<0.3$ against sulfate, nitrate and OC and no significant covariation
with any of five PMF-resolved PM$_{2.5}$ source factors that between them dominated the mass.

\item \textbf{A clean emission-factor dataset for prescribed fire.} Blanco-Alegre et al.\
report PM$_{2.5}$ emission factors of 53.5\,g\,kg$^{-1}$ (\textit{Calluna}, flaming,
MCE 90.6\%) and 44.7\,g\,kg$^{-1}$ (\textit{Genista}, smouldering, MCE 70.8\%), with OC and
EC at 28.1\% and 32.9\% of PM$_{2.5}$ mass. Combustion-regime-resolved EFs of this quality
are scarce for Mediterranean shrubland and directly usable in inventory work.

\item \textbf{Negative result worth recording.} In Lahore's 2024--2025 smog season
(mean PM$_{2.5}$ 96\,$\mu$g/m$^3$, WHO 24\,h guideline exceeded on every monitored day),
routine mask use was \emph{not} associated with lower symptom count --- IRR 0.96
(0.84--1.09) --- while current smoking was, at IRR 1.15 (1.00--1.32). Self-reported mask use
in a cross-sectional sample cannot separate protection from reverse causation, but the null
belongs on record.

\item \textbf{Six of twenty-two records arrived without an abstract} --- four Elsevier and
two RSC deposits, carried with citation and metadata only, flagged as such throughout, and
with no effect size inferred for any of them.
\end{itemize}

\clearpage

% ==================================================================== ANALYTICS
\section*{\sffamily\bfseries\color{Deep}Corpus analytics}
\vspace{-1.5mm}{\color{Amber}\rule{\linewidth}{1.1pt}}\vspace{3mm}

\noindent\includegraphics[width=\linewidth]{\FIGDIR/f1_subtopics.png}
\figcap{Subtopic assignment of the 22 in-scope records. Nine of the ten canonical clusters
are represented --- equalling 14 and 15 August for the widest spread since July, when three
consecutive issues reached all ten. \textit{Reproductive \& developmental} is the single
empty cluster, its ninth absence in twenty-five issues and the first since 13 August.
Assignment is single-label by dominant contribution.}

\vspace{2mm}
\noindent\includegraphics[width=\linewidth]{\FIGDIR/f2_design_endpoint.png}
\figcap{Left --- study architecture. \textit{Metadata only} is the largest single class at
6/22 (27\%), the third-highest share of the watch after 8 August (43\%) and 1 August (32\%);
these are publisher deposits that carried no abstract, not a study design.
Measurement campaigns (4) outnumber observational cohorts (3).
Right --- health endpoint. Six records have no extractable endpoint because no abstract was
deposited, and a further four genuinely carry none (one filtration trial, one emission-factor
study, two aerosol-chemistry campaigns). Respiratory endpoints lead the extractable set at
four.}

\vspace{2mm}
\noindent\includegraphics[width=\linewidth]{\FIGDIR/f3_metric_geography.png}
\figcap{Left --- particle metric. PM$_{2.5}$ alone is the exposure in 10 records and appears
jointly with PM$_{10}$ in a further 5. One record takes ultrafine particle \emph{number} as
its primary metric (a ten-year Athens size-distribution record) and one takes
composition/speciation. Right --- geographic provenance. The \textit{Global / multi-region}
bar aggregates three review syntheses, three records whose study location the abstract does
not state, and two bench studies with no geography. Among located work, China and Europe tie
at four; Sub-Saharan Africa (Addis Ababa) and South Asia (Lahore) contribute one each.}

\vspace{2mm}
\noindent\includegraphics[width=\linewidth]{\FIGDIR/f6_lifecourse.png}
\figcap{Life-course windows addressed. Counts are non-exclusive and a review counts once per
window it explicitly treats. Working-age adulthood dominates (5) --- the CHARLS 45+ wave, the
UK Biobank diabetes cohort, the Korean cirrhosis registry, the Lahore outpatient sample and
the swine-farm workers all sit in that band. The single in-utero entry is a
reviewed determinant of biological-aging trajectory, not a primary prenatal-exposure study ---
today's corpus contains none.}

\vspace{2mm}
\noindent\includegraphics[width=\linewidth]{\FIGDIR/f4_heatmap.png}
\figcap{Subtopic $\times$ study-architecture cross-tabulation. Two structural features stand
out. First, the \textit{Metadata only} column is the most populated in the panel and spans
three different subtopics --- the abstract-free deposit problem is not concentrated in one
literature. Second, the \textit{Sensing, forecasting \& instrumentation} row again contains
no health-endpoint linkage of any kind, and \textit{Exposure assessment \& modelling} carries
exactly one, the Perugia wearable study. That siloing is real today but is not structural:
twelve of the nineteen August issues, this one included, have carried at least one
sensing- or exposure-side record with a human health endpoint.}

\vspace{2mm}
\noindent\includegraphics[width=\linewidth]{\FIGDIR/f5_forest.png}
\figcap{The five quantitative estimates reported in today's corpus, on a common logarithmic
axis. \textbf{The exposure contrasts are not mutually comparable}: Nong et al.\ report
top-versus-bottom quartile hazard ratios over a 12.4-year median follow-up, Park et al.\ a
per-IQR odds ratio on a single day, and the two Ali et al.\ estimates are incidence rate
ratios for \emph{behavioural} covariates (smoking, mask use) measured during --- not
contrasts in --- particulate exposure. Read direction and precision only. Note that the
tightest interval belongs to the estimate closest to the null (Park, 29{,}167 patients) and
that both Ali intervals cross or abut it on an $n$ of 400.}

\clearpage

% ==================================================================== DIGEST
\section*{\sffamily\bfseries\color{Deep}Paper-level digest}
\vspace{-1.5mm}{\color{Amber}\rule{\linewidth}{1.1pt}}\vspace{1mm}

{\fontsize{7.4}{9}\selectfont\color{Slate}Ordered by cluster. Each entry gives the
methodological core and the finding that matters, not an abstract paraphrase. Six records
are marked \textsc{metadata only}: the publisher deposited no abstract, so nothing beyond
title, journal and DOI is asserted for them.\par}

\band{Neuro, cognition \& mental health\hfill 1 record}{Deep}

\paperentry{He et al.\ 2026}
{Plasma proteomic signatures modifying the air pollution--epilepsy association}
{Ecotoxicol Environ Saf}
{Proteome-wide interaction screen: 51{,}428 participants, 2{,}254 Olink-style plasma
proteins, PM$_{2.5}$/PM$_{\mathrm{coarse}}$/NO$_2$/NO from land-use regression, plus a PCA-derived
composite air-pollution index. Logistic models for prevalent and Cox models for incident
epilepsy identify \textbf{CIT--PM$_{2.5}$, PTPRR--NO$_2$ and BCL2L15--API interactions},
enriched for immune, inflammatory and cell-death pathways, with hippocampal-volume and
white-matter-hyperintensity correlates and neuronal/endothelial single-nucleus enrichment.
The threat to inference is multiplicity: $\sim$2{,}254 proteins $\times$ four exposures with
no corrected significance level or interaction magnitude reported in the abstract. A
candidate list, not a mechanism. Relevance here is that it is exposure \emph{misclassification}
in the LUR layer, not protein measurement, that would most easily generate this pattern.}
{10.1016/j.ecoenv.2026.120669}

\band{Cardiovascular \& metabolic\hfill 1 record}{Teal}

\paperentry{Nong et al.\ 2026}
{Long-term air pollution and microvascular complications in diabetes}
{PLoS One}
{UK Biobank, 9{,}671 participants with diabetes, 12.4-year median follow-up, 2{,}104 (21.8\%)
incident microvascular complications. Cox models adjusted for demographic, clinical and
biochemical covariates give \textbf{PM$_{10}$ HR 1.44 (1.27--1.64)} and
\textbf{NO$_2$ HR 1.27 (1.12--1.45)} for the highest versus lowest exposure quartile, stable
across stratified and sensitivity analyses, with early divergence of the cumulative-incidence
curves. Two caveats: PM$_{2.5}$ is listed among the exposures but no PM$_{2.5}$ estimate is
reported in the abstract, and UK Biobank exposure contrasts are narrow, so the quartile
framing hides how small the underlying concentration difference is. Restricted cubic splines
are reported as showing risk beginning at relatively low levels --- the claim that matters
for guideline-adjacent argument, and the one the abstract quantifies least.}
{10.1371/journal.pone.0354711}

\band{Respiratory \& allergic\hfill 3 records}{Sage}

\paperentry{Ali et al.\ 2026}
{Symptom burden among adult outpatients during Lahore's 2024--2025 winter smog}
{Public Health Chall}
{400 outpatients at a tertiary hospital, October 2024 -- January 2025, structured
questionnaire plus city-level air quality (mean AQI 180, mean PM$_{2.5}$ 96\,$\mu$g/m$^3$;
PM$_{2.5}$ and PM$_{10}$ above the WHO 24\,h guideline on every monitored day). 98\% reported
at least one symptom --- breathlessness 74\%, eye irritation 56\%, cough 53\%. Regression on a
composite vulnerability index: \textbf{smoking IRR 1.15 (1.00--1.32)}, \textbf{routine mask
use IRR 0.96 (0.84--1.09), i.e.\ null}. Exposure is assigned at city level to a
health-seeking sample with no unexposed comparison, so no dose--response can be read from
this; its value is documenting symptom prevalence in a South Asian megacity where
hospital-based evidence is thin.}
{10.1002/puh2.70349}

\paperentry{Yang et al.\ 2026a}
{Atmospheric humidity and air pollutants versus self-reported chronic lung disease}
{J Asthma}
{16{,}204 CHARLS adults aged 45+, 2011/2013/2015 waves, six pollutants and six humidity
indices at 1\,km. Each IQR increase in pollutant concentration carried
\textbf{18--45\% higher odds} of self-reported CLD; PM$_{10}$ alone showed an approximately
U-shaped exposure--response, and relative humidity an inverted-U, with a significant
RH$\times$sex interaction. Weighted quantile sum models were positive for both the pollutant
and the humidity mixture. Self-reported outcome and a wide odds range quoted without
per-pollutant intervals limit what can be taken from the point estimates; the reason to read
it is the humidity treatment --- six distinct indices rather than RH alone, which is the same
distinction that separates a well-corrected from a badly-corrected low-cost nephelometer.}
{10.1080/02770903.2026.2721168}

\paperentry{Rice \& Miller 2026}
{Air pollution, microbes and respiratory infection risk in children}
{Pediatr Pulmonol}
{Narrative review joining two literatures usually kept apart: pollutant-driven epithelial
barrier loss, oxidative stress and impaired innate/adaptive immunity, and the developing
airway microbiome as a determinant of colonisation resistance. Argues that dysbiosis
(antibiotics, air pollution, other perturbations) and pollutant exposure are partly
substitutable routes to the same susceptibility phenotype, and notes that exposure to
diverse environmental microbes --- certain fungi included --- may be protective. No
quantitative synthesis and no bias assessment; useful as a framing citation for why an
indoor sensing deployment in a childcare setting should log more than mass.}
{10.1002/ppul.71798}

\band{Mechanistic toxicology\hfill 2 records}{Violet}

\paperentry{Sivakumar \& Kurian 2026}
{Duration-dependent transition from mitochondrial adaptation to persistent cardiac toxicity}
{Environ Toxicol}
{Female Wistar rats, whole-body PM$_{2.5}$ at 250\,$\mu$g/m$^3$, 3\,h/day for 1, 7, 14 or
21 days, with cardiac mitochondrial bioenergetics, redox balance, quality-control signalling
and ex vivo function assessed after each duration and after a 24\,h washout.
\textbf{1--7 days: adaptive and fully reversible. $\geq$14 days: suppressed respiration and
ETC complex activity, mtDNA depletion, impaired biogenesis, intramitochondrial metal
accumulation --- none recovered within the washout.} Concentration was held constant, so
duration alone moved the system across the threshold. The 24\,h washout is short, so
``persistent'' means ``not recovered in 24\,h'', which the authors state plainly; a longer
recovery arm would settle it. The instrumentation implication is direct: a
cumulative-duration-above-threshold metric is testable with continuous logging and invisible
to a 24\,h filter sample.}
{10.1002/tox.70189}

\paperentry{Chen et al.\ 2026a}
{PM$_{2.5}$ suppresses IL-8-mediated epithelial migration in conjunctival cells}
{J Chin Med Assoc}
{Immortalised human conjunctival epithelial line exposed to PM$_{2.5}$ with viability,
apoptosis, scratch-migration and cytokine-array readouts. PM$_{2.5}$ reduced viability
time- and dose-dependently with only mild apoptosis, but markedly impaired migration and
wound closure. Cytokine profiling showed \emph{selective} remodelling rather than a general
inflammatory response, with IL-8 the most strongly suppressed mediator; recombinant IL-8
rescued migration and partially rescued viability, and IL-8 neutralisation impaired migration
at baseline. Single immortalised line, suspension dosing, no dose anchored to an ambient
concentration --- so the mechanism is credible and the dosimetry is not. The ocular surface
remains an under-instrumented exposure site.}
{10.1097/JCMA.0000000000001418}

\band{Occupational \& indoor exposure\hfill 2 records}{Clay}

\paperentry{Xia et al.\ 2026}
{Enhanced filtration for indoor air quality in child-serving facilities}
{Front Pediatr}
{Week-long baseline and follow-up measurements in 33 spaces across 12 facilities spanning
daycare to high school, with staff surveys, as part of a community-based participatory
partnership. Air change rates were low throughout (0.9--1.4\,h$^{-1}$); baseline school-hours
indoor PM$_{2.5}$ 2.2--6.3\,$\mu$g/m$^3$ and PM$_{10}$ 5.7--26.9\,$\mu$g/m$^3$, highest in
standard schools. Enhanced filtration cut PM$_{2.5}$ \textbf{45--54\%} in standard schools and
the nursery and had little effect in childcare houses; a linear mixed model attributes
\textbf{$\sim$51\%} of the PM$_{2.5}$ reduction to the portable air cleaners.
\textbf{Single-zone mass-balance prediction held for regular classrooms and the nursery but
failed for portable buildings, very large rooms and repurposed houses.} That failure boundary
is the most transferable result in the paper: it maps exactly onto where a deployed sensor
earns its keep. No health endpoint was measured.}
{10.3389/fped.2026.1875263}

\paperentry{Park et al.\ 2026b}
{Phenotypic resistance, not community turnover, tracks antimicrobial use in swine-farm bioaerosol}
{Environ Pollut}
{962 isolates across six species, six Korean swine farms stratified by defined daily dose for
animals, four compartments (faecal, aerosol, pig nasal, worker nasal) and two seasons, with
flow-cytometric viability and qPCR-corrected abundance feeding a Monte Carlo QMRA.
\textit{S.\ aureus} isolate-level resistance was higher on high-usage farms; the other four
species and the community composition showed no usage association. \textit{Escherichia}
abundance fell $\sim$four orders of magnitude on aerosolisation while \textit{Staphylococcus}
persisted and peaked in worker nasal samples. Median daily colonisation probability
1.98--8.63$\times$10$^{-4}$; the model was most sensitive to aerosol \emph{genus composition},
not total bioaerosol load. Six farms cannot resolve a farm-level effect and the authors say
so. For sensing: total particle or total bioaerosol load is the wrong exposure variable here.}
{10.1016/j.envpol.2026.128991}

\clearpage

\band{Sensing, forecasting \& instrumentation\hfill 3 records}{Amber}

\paperentry{Yang et al.\ 2026b}
{Dust transport versus local emissions in shaping urban ice-nucleating particles}
{Atmos Chem Phys}
{Month-long winter campaign in Taiyuan, a heavily industrialised Chinese city: immersion-mode
INP concentration and ice-nucleation active-site density alongside size distributions and
chemical composition. $N_{\mathrm{INP}}$ 0.05--13.37\,L$^{-1}$ at and above $-15\,^\circ$C.
During a desert-dust event \textbf{7.47\,L$^{-1}$ (95\% CI 6.64--8.41) versus 0.61\,L$^{-1}$
otherwise}, $n_s$ 1.57$\times$10$^7$ versus 9.52$\times$10$^5$\,m$^{-2}$. During pollution
periods $N_{\mathrm{INP}}$ correlated only weakly with sulfate, nitrate and OC ($|r|<0.3$),
and none of five PMF-resolved PM$_{2.5}$ source factors --- industrial, dust, secondary, coal
and traffic, fireworks --- covaried significantly with it despite dominating the mass. One
city, one winter, one dust event carrying the contrast. Still the cleanest recent
demonstration that a PM$_{2.5}$ mass record does not index a co-located particle property.}
{10.5194/acp-26-11771-2026}

\paperentry{Kalkavouras et al.\ 2026}
{Ten-year trends in ultrafine particle number concentration and size distribution, Athens}
{Atmos Pollut Res}
{\textsc{metadata only} --- no abstract was deposited, so no trend, magnitude or instrument
detail is asserted here. Flagged because a decade-long urban particle-number and
size-distribution record is exactly the reference series that a low-cost optical network
cannot reproduce and must be evaluated against; worth retrieving in full.}
{10.1016/j.apr.2026.103181}

\paperentry{Ayele et al.\ 2026}
{Fine-scale black carbon and TSP composition at a heavy-traffic intersection, Addis Ababa}
{Environ Sci Atmos}
{\textsc{metadata only} --- the deposited abstract is truncated to its opening sentences and
reports no quantities. What is verifiable from it: the study addresses spatiotemporal
variability of equivalent black carbon and the molecular composition of organics in total
suspended particulate at a roadside intersection, and states that Ethiopian urban aerosol
remains poorly characterised. Today's only sub-Saharan record.}
{10.1039/d6ea00100a}

\band{Exposure assessment \& modelling\hfill 4 records}{Amber}

\paperentry{Blanco-Alegre et al.\ 2026}
{Emission factors and particle morphology from prescribed fires in Mediterranean shrubland}
{Atmos Chem Phys}
{Field campaign on \textit{Calluna vulgaris} and \textit{Genista hispanica} burns in
north-west Spain combining gravimetry, thermal--optical transmittance, ion chromatography,
FTIR and SEM. \textit{Calluna} burned flaming (MCE 90.6\%), \textit{Genista} smouldering
(MCE 70.8\%); gas-phase concentrations were higher for \textit{Calluna} but emission factors
for incomplete-combustion products were higher for \textit{Genista}.
\textbf{PM$_{2.5}$ EFs 53.5\,g\,kg$^{-1}$ (\textit{Calluna}) and 44.7\,g\,kg$^{-1}$
(\textit{Genista})}, with OC and EC at 28.1\% and 32.9\% of PM$_{2.5}$ mass and water-soluble
inorganic ions at 6.9\% and 4.5\%. SEM showed abundant tar-ball-like submicron particles and
thickly coated aggregates. Two vegetation types at one site --- the EFs are not a regional
average --- but combustion-regime-resolved and directly usable.}
{10.5194/acp-26-11695-2026}

\paperentry{Duan et al.\ 2026}
{Regime-dependent aerosol pH in ammonia-rich urban Beijing from explainable machine learning}
{Atmos Chem Phys}
{Field observations plus thermodynamic modelling, interrogated with interpretable ML to
apportion pH variability. Temperature dominates, at roughly \textbf{$-0.6$ pH units per
10\,$^\circ$C}. Excess ammonia, nitrate-to-sulfate ratio and PM$_1$ mass are positively
associated with pH, with strongest sensitivity at low values and saturation at high ones.
Relative humidity is the interesting one: its sign \emph{reverses} by regime --- higher RH
increases acidity below $\sim$15\,$^\circ$C, at N/S $>\sim$1.25, or at PM$_1>\sim$50\,$\mu$g/m$^3$,
and decreases it under warm, sulfate-dominant or low-mass conditions. Feature attribution is
association, not causation, and pH is itself a model output rather than a measurement. The
RH regime reversal is a useful caution for anyone applying a single humidity correction
across seasons.}
{10.5194/acp-26-11683-2026}

\paperentry{Mahia et al.\ 2026}
{Street geometry, air quality and real-time physiological response in elderly pedestrians}
{J Environ Manage}
{Fourteen elderly participants walked five routes in Perugia during spring while wearable
sensors logged physiology and portable stations logged microclimate and air quality. Paths
with higher height-to-width ratios (0.92--1.27) and elevated NO$_2$ and PM$_{10}$ (to 21\,ppb
and 18.77\,$\mu$g/m$^3$) significantly degraded air-quality perception
($\beta=-0.273$, $p=0.028$ for NO$_2$; $\beta=-0.312$, $p=0.013$ for PM$_{10}$), and both
pollutants tracked skin-conductance change and self-reported emotional state; higher Green
Visual Index went with more relaxed responses. With $n=14$ over one season the coefficients
carry little weight, and perception and physiology are both plausibly driven by street
geometry directly rather than by the pollutant. Read it for the measurement stack --- paired
personal air-quality and physiological logging at walking cadence --- not the estimates.}
{10.1016/j.jenvman.2026.130681}

\paperentry{Dovrou et al.\ 2026}
{Wintertime wildfires reshape urban atmospheric chemistry: the January 2025 Los Angeles fires}
{Environ Sci Atmos}
{\textsc{metadata only} --- the deposited abstract is truncated before any result. Verifiable
from it: the study concerns how dominant chemical regimes and meteorology govern wildfire
air-quality impacts in urban coastal regions, using the January 2025 Los Angeles fires.
Flagged for retrieval given how much of this watch's recent health signal has come from
source-resolved wildfire PM$_{2.5}$.}
{10.1039/d6ea00096g}

\band{Burden, policy \& mitigation\hfill 1 record}{Coral}

\paperentry{Bhatta et al.\ 2026}
{Outdoor air purification technologies: efficiency, sustainability and health impact}
{Sci Total Environ}
{PRISMA-ScR review, 266 records screened to 28 studies across five technology domains.
Reported bests: smog towers at \textbf{97.8\% PM$_{2.5}$ and 98.4\% PM$_{10}$ removal under
optimal conditions}; solar-chimney hybrids within 2--4\% of dedicated systems while
generating electricity; TiO$_2$-modified cementitious composites potentially carbon-negative
over their lifecycle; a roadside hedge cutting leaf-surface PM$_{2.5}$ magnetic loading 82\%.
The health-economics figure --- 93{,}200 deaths averted in China at a 35\,$\mu$g/m$^3$ indoor
target, positive net monetary benefit in 93\% of cities --- is modelled from \emph{indoor}
air cleaners used as an analogue for outdoor systems, which the authors flag and which should
be read as a bound on plausibility, not an estimate. ``Optimal conditions'' removal
efficiencies are chamber or near-field numbers and say nothing about neighbourhood-scale
concentration. The review's own priority list --- standardised performance metrics, long-term
field validation --- is the honest summary.}
{10.1016/j.scitotenv.2026.182205}

\band{Other clinical endpoints\hfill 5 records}{Slate}

\paperentry{Park et al.\ 2026a}
{Short-term PM$_{2.5}$ and acute hepatic decompensation in cirrhosis}
{Gut Liver}
{Nationwide time-stratified case-crossover on the Korean NHIS database, 2006--2021:
29{,}167 cirrhosis patients presenting to emergency departments with ascites, variceal
bleeding, hepatic encephalopathy, jaundice or hepatorenal syndrome; control days matched on
month and weekday; daily PM$_{2.5}$ assigned on a 1\,km residential grid; conditional logistic
regression adjusted for temperature, humidity, holidays and weekday.
\textbf{aOR 1.02 (1.00--1.04) per IQR increase}, with the top exposure quintile 5\% above the
bottom and \emph{attenuation} at longer averaging windows. No subgroup interaction reached
significance. The design controls time-invariant confounding well, but the lower bound sits
on the null and the effect is small enough that residual exposure misclassification from
grid-assigned concentration could account for it. A large, clean, honestly reported near-null.}
{10.5009/gnl260039}

\paperentry{Poursafa \& Jylh\"av\"a 2026}
{Air pollution and biological aging: mechanisms, biomarkers and reversibility}
{Ageing Res Rev}
{Review organised by biomarker class rather than by outcome, which is what makes it useful.
Evidence is graded as strongest for \emph{system-level composites} --- Klemera--Doubal
biological age, PhenoAge, frailty indices --- and weakest, at meta-analytic level, for
\emph{epigenetic clocks}, which the authors attribute to heterogeneity in exposure assessment,
design and target-tissue cell composition rather than to absent biology. Causal mediation
across independent cohorts supports biological aging as a pathway from exposure to
cardiovascular, neuropsychiatric and renal disease. Pollutants of concern are PM$_{2.5}$,
PM$_{0.1}$, traffic-related pollution and household combustion. Named gaps: household air
pollution and low-income settings are under-represented, and longitudinal multi-omic designs
are scarce. Air-quality and clean-cooking improvements are reported to partially reverse the
acceleration --- the strongest claim in the paper and the one most dependent on the
intervention literature it summarises.}
{10.1016/j.arr.2026.103297}

\paperentry{Wang et al.\ 2026}
{Long-term PM$_{2.5}$, aging hallmarks and age-related disease}
{Atmos Environ}
{\textsc{metadata only} --- no abstract deposited. Title indicates a nationwide prospective
cohort linking long-term PM$_{2.5}$ to aging hallmarks and age-related disease incidence.
Cross-reads directly with Poursafa \& Jylh\"av\"a above, which is the reason to retrieve it
in full rather than wait for indexing.}
{10.1016/j.atmosenv.2026.122313}

\paperentry{Lei et al.\ 2026}
{Hospital admission risk and burden of acute appendicitis attributable to PM$_{2.5}$}
{Atmos Environ}
{\textsc{metadata only} --- no abstract deposited. Title indicates nationwide Chinese
evidence on PM$_{2.5}$ and acute appendicitis admissions with an attributable-burden
component. No effect size, lag structure or exposure model is asserted here.}
{10.1016/j.atmosenv.2026.122311}

\paperentry{Chen et al.\ 2026b}
{Air pollution and sarcopenia: meta-analysis and network toxicology}
{Atmos Environ}
{\textsc{metadata only} --- no abstract deposited. Title indicates a meta-analysis paired
with a network-toxicology arm. Pooled estimate, heterogeneity and included-study count are
all unknown from the deposit; a meta-analysis carried without its numbers is a placeholder,
recorded so it is not missed on a later sweep.}
{10.1016/j.atmosenv.2026.122312}

\clearpage

% ==================================================================== TABLE
\section*{\sffamily\bfseries\color{Deep}Corpus register}
\vspace{-1.5mm}{\color{Amber}\rule{\linewidth}{1.1pt}}\vspace{2.5mm}

{\fontsize{6.7}{8.4}\selectfont
\setlength{\tabcolsep}{2.4pt}
\begin{longtable}{@{}>{\RaggedRight}p{25mm}>{\RaggedRight}p{41mm}>{\RaggedRight}p{27mm}%
>{\RaggedRight}p{23mm}>{\RaggedRight}p{18mm}>{\RaggedRight}p{15mm}@{}}
\toprule
\textsf{\textbf{First author}} & \textsf{\textbf{Focus}} & \textsf{\textbf{Journal}} &
\textsf{\textbf{Design}} & \textsf{\textbf{Metric}} & \textsf{\textbf{Region}}\\
\midrule
\endfirsthead
\toprule
\textsf{\textbf{First author}} & \textsf{\textbf{Focus}} & \textsf{\textbf{Journal}} &
\textsf{\textbf{Design}} & \textsf{\textbf{Metric}} & \textsf{\textbf{Region}}\\
\midrule
\endhead
\multicolumn{6}{@{}l@{}}{\cellcolor{Mist}\textsf{\textbf{\textcolor{Deep}{Neuro / mental health}}}}\\[0.6mm]
\href{https://doi.org/10.1001/jamanetworkopen.2026.25015}{Dales} & Wildfire-sourced PM$_{2.5}$ and ED visits for migraine and other primary headache & \textit{JAMA Netw Open} & Case-crossover & PM$_{2.5}$ (wildfire) & Canada\\
\href{https://doi.org/10.1093/ageing/afag215}{Glover} & Umbrella meta-review of urban environmental exposures and cognitive health & \textit{Age Ageing} & Umbrella review / meta-analysis & PM$_{2.5}$ & Global\\
\href{https://doi.org/10.1016/j.nsa.2026.107019}{Goossens} & In utero air pollution exposure and pre-/postnatal brain development & \textit{Neurosci Appl} & Narrative review & PM$_{2.5}$ / PM$_{10}$ / PM$_{\mathrm{coarse}}$ & Global\\
\href{https://doi.org/10.1016/j.isci.2026.116617}{Ma} & PM$_{2.5}$-induced neurotoxicity and depression: exposure risks and mechanisms & \textit{iScience} & Narrative review & PM$_{2.5}$ + components & Global\\
\href{https://doi.org/10.3760/cma.j.cn112150-20250804-00752}{Pan} & Maternal PM exposure and offspring neurodevelopment: effects and mechanisms & \textit{Zhonghua Yu Fang Yi Xue Za Zhi} & Narrative review & PM (bulk) & China\\
\addlinespace[0.9mm]
\multicolumn{6}{@{}l@{}}{\cellcolor{Mist}\textsf{\textbf{\textcolor{Teal}{Cardiovascular \& metabolic}}}}\\[0.6mm]
\href{https://doi.org/10.3760/cma.j.cn112338-20251120-00843}{Huang} & PM$_{2.5}$ and dyslipidaemia in rural Chinese children; dietary effect modification & \textit{Zhonghua Liu Xing Bing Xue Za Zhi} & Cross-sectional (multistage) & PM$_{2.5}$ & China\\
\href{https://doi.org/10.32394/pe/218966}{Kazmierska} & Air pollution and cardiovascular disease: review of current data & \textit{Przegl Epidemiol} & Narrative review & PM$_{2.5}$ / PM$_{10}$ & Global\\
\href{https://doi.org/10.1007/s00125-026-06814-2}{Xi} & Air pollution, built environment and temperature vs. prediabetes phenotypes & \textit{Diabetologia} & Prospective cohort & PM$_{2.5}$ & Germany\\
\href{https://doi.org/10.1161/JAHA.125.046917}{Yang} & Short-term PM$_{2.5}$ and hospitalisation/ED use in atrial fibrillation & \textit{J Am Heart Assoc} & Case-crossover & PM$_{2.5}$ (wildfire-influenced) & USA\\
\addlinespace[0.9mm]
\multicolumn{6}{@{}l@{}}{\cellcolor{Mist}\textsf{\textbf{\textcolor{Sky}{Reproductive \& developmental}}}}\\[0.6mm]
\href{https://doi.org/10.1016/j.jes.2025.10.043}{Chen} & Prenatal air pollution, fetal growth trajectory phenotypes and child neurodevelopment & \textit{J Environ Sci (China)} & Prospective cohort & PM$_{2.5}$ / PM$_{10}$ & China\\
\href{https://doi.org/10.1016/j.ecoenv.2026.120566}{Huang} & Joint prenatal bisphenol + air-pollutant mixture exposure and preschool IQ & \textit{Ecotoxicol Environ Saf} & Prospective cohort (mixtures) & PM$_{2.5}$ / PM$_{10}$ & China\\
\href{https://doi.org/10.1007/s40572-026-00552-8}{Peterson} & Ambient air pollution, PAHs and pubertal development: critical review & \textit{Curr Environ Health Rep} & Critical review & PM$_{2.5}$ components & Global\\
\href{https://doi.org/10.1016/j.envres.2026.125313}{Xie} & Long-term PM$_{2.5}$ and timing of pubertal development in boys and girls & \textit{Environ Res} & Prospective cohort & PM$_{2.5}$ & China\\
\addlinespace[0.9mm]
\multicolumn{6}{@{}l@{}}{\cellcolor{Mist}\textsf{\textbf{\textcolor{Sage}{Respiratory \& allergic}}}}\\[0.6mm]
\href{https://doi.org/10.1007/s00484-026-03289-z}{Bayar Muluk} & Climate change and otolaryngology: airway disease impacts and emission reduction & \textit{Int J Biometeorol} & Narrative review & PM$_{2.5}$ & Global\\
\href{https://doi.org/10.4103/mgr.MEDGASRES-D-25-00164}{Liao} & Short-term air pollution, respiratory symptoms and lung function in asthma & \textit{Med Gas Res} & Panel study & PM$_{2.5}$ / PM$_{10}$ & China\\
\href{https://doi.org/10.3390/epidemiologia7040097}{Sanchez de Prada} & Ambient PM and asthma-coded admissions in children 0-3 y, Spain & \textit{Epidemiologia} & Ecological panel & PM$_{2.5}$ / PM$_{10}$ & Spain\\
\href{https://doi.org/10.1016/j.jhazmat.2026.143076}{Yang} & Allergenic vs. toxic risk of airborne seed hairs compared with ambient PM$_{2.5}$ & \textit{J Hazard Mater} & Field measurement & PM$_{2.5}$ + bioaerosol & China\\
\addlinespace[0.9mm]
\multicolumn{6}{@{}l@{}}{\cellcolor{Mist}\textsf{\textbf{\textcolor{Violet}{Mechanistic toxicology}}}}\\[0.6mm]
\href{https://doi.org/10.1016/j.jes.2025.12.074}{Ding} & PM$_{2.5}$ promotes vascular calcification via miR-27b-5p/HES1 signalling & \textit{J Environ Sci (China)} & In vitro + organoid + murine & PM$_{2.5}$ & China\\
\href{https://doi.org/10.1007/s00360-026-01698-5}{Maraschi} & Osmoregulatory disruption and metal bioaccumulation from settleable atmospheric PM & \textit{J Comp Physiol B} & Controlled exposure (ecotox) & Settleable APM & Brazil\\
\href{https://doi.org/10.1016/j.jes.2025.11.018}{Wu} & Macrophage NLRP3 inflammasome in real-ambient PM-induced hepatic glucose dysregulation & \textit{J Environ Sci (China)} & Real-ambient murine + co-culture & Real-ambient PM & China\\
\addlinespace[0.9mm]
\multicolumn{6}{@{}l@{}}{\cellcolor{Mist}\textsf{\textbf{\textcolor{Clay}{Occupational \& indoor}}}}\\[0.6mm]
\href{https://doi.org/10.3389/fmedt.2026.1796828}{Chen} & Nested suction instrument to cut intraoperative surgical-smoke PM$_{2.5}$ & \textit{Front Med Technol} & Bench / device experiment & PM$_{2.5}$ (surgical smoke) & China\\
\href{https://doi.org/10.1093/annweh/wxag055}{Cole} & Determinants of respirable crystalline silica exposure in Australian tunnelling & \textit{Ann Work Expo Health} & Observational exposure study & Respirable dust / RCS & Australia\\
\href{https://doi.org/10.1016/j.jhazmat.2026.143081}{Ding} & Stage-resolved hair Zn-isotope dynamics as a firefighting combustion-exposure tracer & \textit{J Hazard Mater} & Repeated-measures biomarker & Combustion-derived PM & China\\
\href{https://doi.org/10.3390/jox16040119}{Ramsden} & Xenobiotic hazards in aircraft cabin air & \textit{J Xenobiot} & Review + risk economics & Ultrafine particles & Global\\
\addlinespace[0.9mm]
\multicolumn{6}{@{}l@{}}{\cellcolor{Mist}\textsf{\textbf{\textcolor{Amber}{Sensing, forecasting \& instrumentation}}}}\\[0.6mm]
\href{https://doi.org/10.1016/j.jenvman.2026.130515}{Shi} & Adaptive spatiotemporal graph transformer for multi-site multi-step PM$_{2.5}$ forecasting & \textit{J Environ Manage} & Deep-learning model & PM$_{2.5}$ & China / India\\
\href{https://doi.org/10.3390/bios16070380}{Tsuchiya} & Seasonal and indoor-outdoor characteristics of fluorescent aerosol particles in offices & \textit{Biosensors} & Real-time sensor campaign & Bioaerosol / coarse & Japan\\
\addlinespace[0.9mm]
\multicolumn{6}{@{}l@{}}{\cellcolor{Mist}\textsf{\textbf{\textcolor{Amber}{Exposure assessment \& modelling}}}}\\[0.6mm]
\href{https://doi.org/10.1029/2025GH001621}{Solorzano-Araque} & Multiscale exposure modelling: LOTOS-EUROS + mobility surveys + particle toxicity & \textit{GeoHealth} & CTM + mobility integration & PM$_{2.5}$ / PM$_{10}$ & Colombia\\
\addlinespace[0.9mm]
\multicolumn{6}{@{}l@{}}{\cellcolor{Mist}\textsf{\textbf{\textcolor{Coral}{Burden, policy \& mitigation}}}}\\[0.6mm]
\href{https://doi.org/10.1002/hsr2.72886}{Karamzad} & Ischaemic heart disease burden in MENA 1990-2021; PM attributable fraction 28.2\% & \textit{Health Sci Rep} & GBD secondary analysis & Ambient PM & MENA\\
\href{https://doi.org/10.1371/journal.pone.0347989}{Rees} & HEPA filtration to reduce respiratory infection in care homes & \textit{PLoS One} & RCT process evaluation & Filtered PM & UK\\
\href{https://doi.org/10.1016/j.jes.2026.01.009}{Shan} & Policy-driven emission reductions and post-pandemic rebound, high-resolution inventories & \textit{J Environ Sci (China)} & Emission inventory & PM emissions & China\\
\addlinespace[0.9mm]
\multicolumn{6}{@{}l@{}}{\cellcolor{Mist}\textsf{\textbf{\textcolor{Slate}{Other clinical endpoints}}}}\\[0.6mm]
\href{https://doi.org/10.3390/dj14070440}{Lupita} & PM$_{2.5}$, preventive-care capacity and caries/edentulism burden across the EU & \textit{Dent J} & Ecological panel & PM$_{2.5}$ & EU-27\\
\href{https://doi.org/10.1016/j.jdin.2026.06.008}{Park} & Long-term PM exposure and risk of androgenetic alopecia & \textit{JAAD Int} & Nationwide cohort & PM$_{2.5}$ / PM$_{10}$ & South Korea\\
\href{https://doi.org/10.3171/2026.3.PEDS2615}{Sun} & Environmental exposure and area deprivation vs. paediatric hydrocephalus outcomes & \textit{J Neurosurg Pediatr} & Registry cohort + geospatial & PM$_{2.5}$ (satellite-derived) & USA\\
\addlinespace[0.9mm]
\bottomrule
\end{longtable}}


\vspace{4mm}
\begin{keybox}
{\sffamily\bfseries\footnotesize\color{Deep}Provenance and method}\\[1.2mm]
{\fontsize{7.2}{8.8}\selectfont
Entry window \textbf{19 August 2026}, contiguous with the 18 August issue. Sources queried
on 19 August: \textbf{PubMed} E-utilities and the PubMed connector, run as two separate
queries against \texttt{[Date - Entry]} --- a health axis and a monitoring/sensing axis
(18 unique PMIDs between them, 14 local and 4 connector-only);
\textbf{Crossref by ISSN} across the instrumentation and atmospheric journals PubMed does not
index (17 deposits over 18 journals); \textbf{Europe PMC} \texttt{CREATION\_DATE}
(0 records --- a topic-free control query returned only 0 records of \emph{any} kind created
that day against 4 on 18 August and 2{,}794 on 17 August, so the index had not ingested the
day rather than the leg being broken); \textbf{arXiv} relevance query with a three-day
recency screen (0 in window, newest posting predates it); \textbf{Consensus}
(no in-window record, tenth consecutive zero-carry); \textbf{OpenAlex} HTTP 429, seventeenth
consecutive failure --- that leg remains dead. ClinicalTrials.gov returned no new PM-related
registration; \texttt{state/trials.json} is unchanged and is surfaced in the weekly rollup.
Deduplication keys on PMID, then lowercased DOI, then a title/author/year signature, checked
against the running \texttt{seen.json} index before any record is summarised; no duplicate
was found today. \textbf{13 records were rejected} with a logged reason --- greenhouse-gas and
carbon-accounting work, cloud and gas-phase atmospheric chemistry, an e-cigarette aerosol
study, and a commentary with no abstract --- all in \texttt{state/rejected.jsonl}.
\textbf{Six of the 22 in-scope records carry no abstract} because the publisher deposited
none; they are marked \textsc{metadata only} in the digest and as ``Metadata only'' in the
register and figures, and nothing beyond title, journal and DOI is asserted for them.
Subtopic, design, particle-metric and geography labels are assigned from title, abstract and
metadata; assignment is single-label by dominant contribution and is a judgement, not a
controlled vocabulary. Every DOI was machine-checked against its PubMed record and resolved
before the build (0 fail, 0 warn). Content derived from publisher and PubMed metadata;
abstracts remain the copyright of their respective publishers.\par}
\end{keybox}

\end{document}
